\documentclass[]{book}

%These tell TeX which packages to use.
\usepackage{array,epsfig}
\usepackage{amsmath}
\usepackage{amsfonts}
\usepackage{amssymb}
\usepackage{amsxtra}
\usepackage{amsthm}
\usepackage{mathrsfs}
\usepackage{color}
\usepackage{mathtools}
\usepackage{geometry}


%Here I define some theorem styles and shortcut commands for symbols I use often
\theoremstyle{definition}
\newtheorem{defn}{Definition}
\newtheorem{thm}{Theorem}
\newtheorem{cor}{Corollary}
\newtheorem*{rmk}{Remark}
\newtheorem{lem}{Lemma}
\newtheorem*{joke}{Joke}
\newtheorem{ex}{Example}
\newtheorem*{soln}{Solution}
\newtheorem{prop}{Proposition}

\newcommand{\lra}{\longrightarrow}
\newcommand{\ra}{\rightarrow}
\newcommand{\surj}{\twoheadrightarrow}
\newcommand{\graph}{\mathrm{graph}}
\newcommand{\bb}[1]{\mathbb{#1}}
\newcommand{\Z}{\bb{Z}}
\newcommand{\Q}{\bb{Q}}
\newcommand{\R}{\bb{R}}
\newcommand{\C}{\bb{C}}
\newcommand{\N}{\bb{N}}
\newcommand{\M}{\mathbf{M}}
\newcommand{\m}{\mathbf{m}}
\newcommand{\MM}{\mathscr{M}}
\newcommand{\HH}{\mathscr{H}}
\newcommand{\Om}{\Omega}
\newcommand{\Ho}{\in\HH(\Om)}
\newcommand{\bd}{\partial}
\newcommand{\del}{\partial}
\newcommand{\bardel}{\overline\partial}
\newcommand{\textdf}[1]{\textbf{\textsf{#1}}\index{#1}}
\newcommand{\img}{\mathrm{img}}
\newcommand{\ip}[2]{\left\langle{#1},{#2}\right\rangle}
\newcommand{\inter}[1]{\mathrm{int}{#1}}
\newcommand{\exter}[1]{\mathrm{ext}{#1}}
\newcommand{\cl}[1]{\mathrm{cl}{#1}}
\newcommand{\ds}{\displaystyle}
\newcommand{\vol}{\mathrm{vol}}
\newcommand{\cnt}{\mathrm{ct}}
\newcommand{\osc}{\mathrm{osc}}
\newcommand{\LL}{\mathbf{L}}
\newcommand{\UU}{\mathbf{U}}
\newcommand{\support}{\mathrm{support}}
\newcommand{\AND}{\;\wedge\;}
\newcommand{\OR}{\;\vee\;}
\newcommand{\Oset}{\varnothing}
\newcommand{\st}{\ni}
\newcommand{\wh}{\widehat}
\newcommand{\F}{\mathcal{F}}
\newcommand{\Fp}{\mathcal{F}_+}
\newcommand{\Tr}{\operatorname{Tr}}
\newcommand{\inner}[2]{\langle #1, #2 \rangle}
\newcommand{\adag}[1]{a^{*}_{#1}}
\newcommand{\an}[1]{a_{#1}}
\newcommand{\LN}{\mathcal{L}_N}
\newcommand{\HN}{\mathcal{H}_N}

%Pagination stuff.
\setlength{\topmargin}{-.3 in}
\setlength{\oddsidemargin}{0in}
\setlength{\evensidemargin}{0in}
\setlength{\textheight}{9.in}
\setlength{\textwidth}{6.5in}
\pagestyle{empty}



\begin{document}


\begin{center}
{\Large Esercizi MMMQ}\\
%textbf{NAME}\\ %You should put your name here
%Due: DATE %You should write the date here.
\end{center}

\vspace{0.2 cm}




\begin{enumerate}
\item Sia $\hat{x}: D(\hat{x}) \to L^2(\mathbb{R})$ definito da $\hat{x}(\psi)(x) = x \psi (x)$ per ogni $\psi \in D(\hat{x}) = \{\psi \in L^2(\mathbb{R}): \int_{\mathbb{R}}x^2|\psi(x)|^2dx < \infty\}$. Si dimostri che $D(\hat{x})$ è denso in $L^2(\mathbb{R})$.

\item Dimostrare che $C^\infty(\mathbb{R})$ è denso in $L^2(\mathbb{R})$.

\item Sia $f: \mathbb{R} \to \mathbb{C}$, si condiseri $(A_f \psi)(x) = f(x)\psi(x)$ per ogni $\psi in D(A) = \{\psi \in L^2(\mathbb{R}): f \psi \in L^2(\mathbb{R})\}$. Si dimostri che $D(A)$ è denso in $L^2(\mathbb{R})$.

\item Siano $D(A) = C^{\infty}_0(\mathbb{R})$, $H = L^2(\mathbb{R})$ e $Af = \sum_{n \in \mathbb{N}}f(n)e_n$, dove $\{e_n: n \in \mathbb{N}\}$ è un sistema ortonormale completo. Dimostrare che $D(A^\star) = \{0\}$.

\item Siano $H = L^2(\mathbb{R})$, $D(A) = \{f \in L^2(\mathbb{R}): \int_\mathbb{R} |xf(x)|dx < \infty\}$ e $(Af)(x) = xf(x)$. Dimostrare che $\sigma(A) = \mathbb{R}, \sigma_p(A) = \emptyset, \rho(A) = \mathbb{C}\setminus \mathbb{R}$.

\item Sia $X = C([0, 1])$ con $||f|| = \sup_{x \in [0, 1]}|f(x)|$. Consideriamo gli operatori $A_i : D_i \subset X \to X$ definiti da $(A_i f)(x) = f'(x)$ per $i = 1, 2, 3, 4$ dove: $D_1 = C^1([0, 1]), D_2 = \{f \in X: f(0) = 0\}, D_3 = \{f \in X: f(0) = f(1)\}$ e $D_4 = \{f \in X: f(0) = f(1) = 0\}$.
Dimostrare che:
\begin{enumerate}
    \item $\sigma(A_1) = \sigma_p(A_1) = \mathbb{C}$ e $\rho(A_1) = \emptyset$;
    \item $\sigma(A_1) = \sigma_p(A_1) = \emptyset$ e $\rho(A_2) = \mathbb{C}$;
    \item $\sigma(A_3) = \sigma_p(A_3) = 2\pi i \mathbb{Z}$;
    \item $\sigma(A_4) = \mathbb{C}, \sigma_p(A_4) = \rho(A_4) = \emptyset.$
\end{enumerate}

\item Sia $u \in H^m(\mathbb{R}^n)$ con $2m \geq n$. Si usi la trasformata di Fourier per dimostrare che $u \in L^{\infty}(\mathbb{R}^n)$ e 
\[
||u||_{L^\infty} \leq C_{n, m}||u||_{H^m},
\]
dove $C_{n,m}$ non dipende da $u$.

\item Siano $A$ un operatore autoaggiunto, $B_n = in(A+in)^{-1}$, $A_n = B_n A B_{-n}$ con $n \in \mathbb{Z}$. Dimostrare che
\[
U_n(t) = e^{iA_n t} = \sum_{k \geq 0} \frac{(-itA_n)^k}{k!}
\]
è un gruppo unitario fortemente continuo.

\item Sia $H_0 = -\Delta:H^2(\mathbb{R}^3) \to L^2(\mathbb{R}^3)$. Mostrare che per $\varphi \in L^2(\mathbb{R}^3)$, $k > 0$, vale che:
\[
((H_0 + k^2)^{-1}\varphi)(x) = \frac{1}{4\pi}\int_{\mathbb{R}^3} \frac{e^{-k|x-y|}}{|x-y|}\varphi(y)dy.
\]Hint: usa la trasformata di Fourier e il teorema dei residui.

\item Sia $B \in \mathcal{L}(H)$. Dimostrare che:
\begin{enumerate}
    \item Se $B \geq 0$ allora $||B|| = \sup_{||\varphi|| = 1}\langle \varphi, B\varphi \rangle$;
    \item $||B^\star B|| = ||B||^2.$
\end{enumerate}

\item Sia \(V \in L^{\infty} + L^2\) a valori reali. Dimostrare che \(\forall \epsilon > 0\) \(\exists V_2 \in L^2(\mathbb{R}^d), \exists V_{\infty} \in L^{\infty}(\mathbb{R}^d)\) 
tali che \(V = V_2 + V_{\infty}\) con \(||V_2||_2 < \epsilon\).

\item Siano \(f, g \in L^2(\mathbb{R}^3)\), si dimostri che;
\begin{enumerate}
    \item \(W(f)\) è unitario e \(W^*(f) = W(f)^{-1} = W(-f)\);
    \item \(W^*(f)a_xW(f) = a_x + f(x)\) e \(W^*(f)a_x^*W(f)= a_x^* + \overline{f}(x)\);
    \item \(\langle W(f)\Omega, \mathcal{N}W(f)\Omega \rangle = ||f||_2^2\).
\end{enumerate}
Dove \(W(f) := e^{a^{*}(f)-a(f)}\) per \(f \in L^2(\mathbb{R}^3)\) e $a^{*}(f)$ e $a(f)$ sono gli operatori di creazione e annichilazione
sullo spazio di Fock $\F$. 


\item Dimostrare che se \(\gamma_{\psi}\) è la matrice di densitò ridotta a una particella, allora \(\Tr \gamma_{\psi} = \langle \psi, \mathcal{N}\psi\rangle \).
\item Consideriamo l'evoluzione temporale di Stati Coerenti:
\[
\begin{cases}
i\,\partial_t\, \psi_{N,t} = \HN\, \psi_{N,t} \\[2pt]
\psi_{N,0} = W(\sqrt{N}\,\varphi)\,\Om,
\end{cases}
\qquad \varphi \in H^1(\mathbb{R}^3),\ \|\varphi\|_2 = 1.
\]
Definiamo
\[
U_N(t,s) = W^{*}(\sqrt{N}\,\varphi_t)\; e^{-i\HN(t-s)}\; W(\sqrt{N}\,\varphi_s),
\]
che soddisfa
\[
i\,\partial_t\, U_N(t,s) = \LN(t)\, U_N(t,s),
\qquad U_N(s,s) = \mathbb{1},
\]
con generatore
\[
\LN(t)
= \Big( i\,\partial_t\, W^{*}(\sqrt{N}\,\varphi_t)\Big) W(\sqrt{N}\,\varphi_t)
+ W^{*}(\sqrt{N}\,\varphi_t)\, \HN\, W(\sqrt{N}\,\varphi_t).
\]
Dimostrare che il generatore $\LN(t)$ si scrive esplicitamente come
\begin{align*}
\LN(t) &= \HN
+ \int dx \,\big(V * |\varphi_t|^2\big)(x)\, \adag{x}\, \an{x} \\
&\quad + \int dx\,dy\; V(x-y)\, \overline{\varphi_t(x)}\, \varphi_t(y)\, \adag{y}\, \an{x} \\
&\quad + \frac{1}{2} \int dx\,dy\; V(x-y)
   \Big( \varphi_t(x)\, \varphi_t(y)\, \adag{x}\, \adag{y}
   + \overline{\varphi_t(x)}\, \overline{\varphi_t(y)}\, \an{x}\, \an{y} \Big) \\
&\quad + \frac{1}{\sqrt{N}} \int dx\,dy\; V(x-y)\, \adag{x}
   \big( \varphi_t(y)\, \adag{y} + \overline{\varphi_t(y)}\, \an{y} \big)\, \an{x}.
\end{align*}
\item 
\begin{enumerate}
\item \(A\) è un proiettore di rango \(1\), allora \(||A||_{\mathcal{I}_1} = 2||A||_{\mathcal{I}_2}\).
\item Usare il punto precedente per dimostrare che:
\[
\Tr\big| \gamma_{N,t} - |\varphi_t\rangle \langle\varphi_t| \big| = 2\,\big\| \gamma_{N,t} - |\varphi_t\rangle \langle\varphi_t| \big\|_{\mathcal{I}_2},
\]
dove \(\gamma_{N, t}\) è la matrice densità ridotta associata allo stato \(\psi_{N,t}\) definito nell'esercizio precedente.
\end{enumerate}
\end{enumerate}
\end{document}
